%% ============================================================
%% AULA 6 — A rede que decora: sobreajuste, regularização e
%%          generalização
%% GBC073 — Inteligência Computacional (FACOM/UFU)
%%
%% Esqueleto com esboço de conteúdo — a desenvolver.
%% Ementa (ficha SEI 5116656): item 2 — projeto de redes (topologia,
%% parâmetros).
%%
%% Compilar:  latexmk -xelatex aula06.tex   (duas passadas!)
%% ============================================================
\documentclass[aspectratio=169, 11pt]{beamer}

\makeatletter
\def\input@path{{./}{../}}
\makeatother

\usetheme{InteligenciaComputacional}   % ANTES do babel: fontspec vem primeiro
\usepackage{babel}
\babelprovide[import, main]{brazilian}

\title[Aula 6 — Generalização]{Aula 6: A rede que decora — sobreajuste, regularização e generalização}
\subtitle[GBC073 — Inteligência Computacional]{Viés e variância \textbullet\ validação \textbullet\ dropout \textbullet\ weight decay}
\author[Prof.\ Marcelo Albertini]{Prof.\ Marcelo Keese Albertini}
\institute{GBC073 — Inteligência Computacional \textbullet\ Universidade Federal de Uberlândia}
\date{2026}

\begin{document}

% ------------------------------------------------------------ capa
\begin{frame}[plain, noframenumbering]
  \titlepage
\end{frame}

% ------------------------------------------------------------ roteiro
\begin{frame}{Roteiro da aula}
  \begin{itemize}\setlength{\itemsep}{0.5ex}
    \item<1-> \textbf{O sobreajuste e o dilema viés-variância}
    \item<2-> \textbf{Validação, early stopping e conjunto de teste}
    \item<3-> \textbf{Regularização explícita:} weight decay e dropout
    \item<4-> \textbf{Normalização e aumento de dados}
  \end{itemize}
\end{frame}

% ------------------------------------------------------------
\section{O sobreajuste e o dilema viés-variância}

\begin{frame}{O sobreajuste e o dilema viés-variância}
  \begin{itemize}\setlength{\itemsep}{0.4ex}
    \item Sobreajuste (\emph{overfitting}): decorar o conjunto de treino e
          falhar no que não viu.
    \item Erro de generalização = viés (modelo simples demais) + variância
          (modelo sensível demais ao sorteio dos dados) + ruído.
    \item Redes grandes têm variância alta — daí a fama de ``decorar''.
    \item Mas redes grandes \emph{com} regularização costumam generalizar
          melhor (contra-intuição importante).
  \end{itemize}
\end{frame}

% ------------------------------------------------------------
\section{Validação, early stopping e conjunto de teste}

\begin{frame}{Validação, early stopping e conjunto de teste}
  \begin{itemize}\setlength{\itemsep}{0.4ex}
    \item Três partições: treino (ajusta), validação (decide), teste (relata
          — só uma vez).
    \item \emph{Early stopping}: parar quando o erro de validação para de cair.
    \item Escolher hiperparâmetros pelo teste vaza o teste — vício sutil.
    \item Curvas de aprendizado como ferramenta de diagnóstico.
  \end{itemize}
\end{frame}

% ------------------------------------------------------------
\section{Regularização explícita}

\begin{frame}{Regularização explícita (weight decay, dropout)}
  \begin{itemize}\setlength{\itemsep}{0.4ex}
    \item Weight decay ($L_2$ nos pesos): penaliza pesos grandes, encolhe para
          soluções simples.
    \item Dropout: sorteia neurônios desligados a cada passo — a rede aprende a
          ser redundante.
    \item Dropout como \emph{ensemble} barato: média implícita de muitas redes.
    \item Em teste, nada é desligado (ou se re-escala a ativação).
  \end{itemize}
  \begin{ideiabox}
    Regularizar não é castigar a rede — é impedir que ela memorize o que não
    deveria existir como regra.
  \end{ideiabox}
\end{frame}

% ------------------------------------------------------------
\section{Normalização e aumento de dados}

\begin{frame}{Normalização e aumento de dados}
  \begin{itemize}\setlength{\itemsep}{0.4ex}
    \item Normalização de lote (\emph{batch norm}): estabiliza a distribuição
          das ativações — treina mais rápido e regulariza de brinde.
    \item Aumento de dados (\emph{data augmentation}): fabricar exemplos
          plausíveis (rotações, ruído, recortes).
    \item A regularização mais barata costuma ser mais dados — reais ou
          sintéticos.
  \end{itemize}
\end{frame}

\end{document}
